
\documentclass{article}
\usepackage{colortbl}
\usepackage{makecell}
\usepackage{multirow}
\usepackage{supertabular}

\begin{document}

\newcounter{utterance}

\centering \large Interaction Transcript for game `dond', experiment `semi\_en', episode 9 with Qwen3{-}32B{-}t0.0.
\vspace{24pt}

{ \footnotesize  \setcounter{utterance}{1}
\setlength{\tabcolsep}{0pt}
\begin{supertabular}{c@{$\;$}|p{.15\linewidth}@{}p{.15\linewidth}p{.15\linewidth}p{.15\linewidth}p{.15\linewidth}p{.15\linewidth}}
    \# & $\;$A & \multicolumn{4}{c}{Game Master} & $\;\:$B\\
    \hline

    \theutterance \stepcounter{utterance}  
    & & \multicolumn{4}{p{0.6\linewidth}}{
        \cellcolor[rgb]{0.9,0.9,0.9}{
            \makecell[{{p{\linewidth}}}]{
                \texttt{\tiny{[P1$\langle$GM]}}
                \texttt{You are playing a negotiation game in which you have to agree on how to divide a set of items among you and another player.} \\
\\ 
\texttt{Rules:} \\
\texttt{(a) You and the other player are given a set of items. Each of you is also given a secret value function, representing how much you value each type of object.} \\
\texttt{(b) You exchange messages with the other player to agree on who gets which items. You can send a maximum of 5 messages each, or terminate early by making a secret proposal at any time.} \\
\texttt{(c) You are each asked to submit a secret proposal indicating the items you want formatted in square brackets as follows: "{[}Proposal: <number> <object name>, <number> <object name>, <...>{]}"} \\
\texttt{(d) If your proposals are complementary, i.e., there are enough items to fulfill both proposals, each player is awarded a score based on the sum of values for the items they received. Otherwise, both of you get zero points.} \\
\texttt{(e) Your goal is to maximize the score you receive.} \\
\\ 
\texttt{Let us start.} \\
\\ 
\texttt{The set of available items is:} \\
\\ 
\texttt{3 candies, 2 fruits, 1 hat} \\
\\ 
\texttt{Your secret value function is:} \\
\\ 
\texttt{candy: 2, fruit: 2, hat: 0} \\
\\ 
\texttt{IMPORTANT: Your messages, unless it is the secret proposal, are directly transmitted to the other player, so do not include any response to the rules or text announcing your message. To make a secret proposal, use the indicated format. Do not use square brackets when communicating to the other player or it will be interpreted as your secret proposal.} \\
            }
        }
    }
    & \\ \\

    \theutterance \stepcounter{utterance}  
    & \multicolumn{4}{p{0.6\linewidth}}{
        \cellcolor[rgb]{0.9,0.9,0.9}{
            \makecell[{{p{\linewidth}}}]{
                \texttt{\tiny{[P1$\rangle$GM]}}
                \texttt{<think>} \\
\texttt{Okay, let's see. I need to negotiate the division of 3 candies, 2 fruits, and 1 hat with another player. My value function is candy: 2, fruit: 2, hat: 0. So the hat is worthless to me, but maybe the other person values it. My goal is to maximize my score, which is the sum of the values of the items I get.} \\
\\ 
\texttt{First, I should figure out what the other player's value function is. But since it's secret, I don't know. However, maybe I can infer it through negotiation. The hat is the only item I don't value, so if the other player values the hat, maybe I can offer them the hat in exchange for candies and fruits.} \\
\\ 
\texttt{Since I can send up to 5 messages, I need to plan my approach. Let's start by making an initial offer. Maybe I can propose taking all the candies and fruits, leaving the hat for them. But there are 3 candies and 2 fruits. If I take all, that's 3*2 + 2*2 = 6 + 4 = 10 points for me. But the other player might not agree if they also want some candies or fruits. Alternatively, maybe a split where I take more candies and they take more fruits, depending on their preferences.} \\
\\ 
\texttt{Wait, the other player might value the hat. Let me think. If I offer them the hat in exchange for some candies or fruits, that could work. For example, if I take 2 candies and 2 fruits, and give them the hat and 1 candy. Then my total would be 2*2 + 2*2 = 8, and they get 1 candy (value depends on their function) plus the hat. But I don't know their values. Maybe I should ask them what they value more.} \\
\\ 
\texttt{But since I can't ask directly, maybe I can make offers and see their response. Let's start with a message proposing a division. Let me try to suggest taking 2 candies and 2 fruits, leaving 1 candy and the hat for them. But I need to check if that's possible. The total items are 3 candies and 2 fruits. If I take 2 candies and 2 fruits, that leaves 1 candy and 0 fruits and 1 hat. So the other player would get 1 candy and 1 hat.} \\
            }
        }
    }
    & & \\ \\

    \theutterance \stepcounter{utterance}  
    & & & \multicolumn{2}{p{0.3\linewidth}}{
        \cellcolor[rgb]{0.9,0.9,0.9}{
            \makecell[{{p{\linewidth}}}]{
                \texttt{\tiny{[GM$|$GM]}}
                \texttt{continue} \\
            }
        }
    }
    & & \\ \\

    \theutterance \stepcounter{utterance}  
    & & & \multicolumn{2}{p{0.3\linewidth}}{
        \cellcolor[rgb]{0.9,0.9,0.9}{
            \makecell[{{p{\linewidth}}}]{
                \texttt{\tiny{[GM$|$GM]}}
                \texttt{<think>} \\
\texttt{Okay, let's see. I need to negotiate the division of 3 candies, 2 fruits, and 1 hat with another player. My value function is candy: 2, fruit: 2, hat: 0. So the hat is worthless to me, but maybe the other person values it. My goal is to maximize my score, which is the sum of the values of the items I get.} \\
\\ 
\texttt{First, I should figure out what the other player's value function is. But since it's secret, I don't know. However, maybe I can infer it through negotiation. The hat is the only item I don't value, so if the other player values the hat, maybe I can offer them the hat in exchange for candies and fruits.} \\
\\ 
\texttt{Since I can send up to 5 messages, I need to plan my approach. Let's start by making an initial offer. Maybe I can propose taking all the candies and fruits, leaving the hat for them. But there are 3 candies and 2 fruits. If I take all, that's 3*2 + 2*2 = 6 + 4 = 10 points for me. But the other player might not agree if they also want some candies or fruits. Alternatively, maybe a split where I take more candies and they take more fruits, depending on their preferences.} \\
\\ 
\texttt{Wait, the other player might value the hat. Let me think. If I offer them the hat in exchange for some candies or fruits, that could work. For example, if I take 2 candies and 2 fruits, and give them the hat and 1 candy. Then my total would be 2*2 + 2*2 = 8, and they get 1 candy (value depends on their function) plus the hat. But I don't know their values. Maybe I should ask them what they value more.} \\
\\ 
\texttt{But since I can't ask directly, maybe I can make offers and see their response. Let's start with a message proposing a division. Let me try to suggest taking 2 candies and 2 fruits, leaving 1 candy and the hat for them. But I need to check if that's possible. The total items are 3 candies and 2 fruits. If I take 2 candies and 2 fruits, that leaves 1 candy and 0 fruits and 1 hat. So the other player would get 1 candy and 1 hat.} \\
            }
        }
    }
    & & \\ \\

    \theutterance \stepcounter{utterance}  
    & & \multicolumn{4}{p{0.6\linewidth}}{
        \cellcolor[rgb]{0.9,0.9,0.9}{
            \makecell[{{p{\linewidth}}}]{
                \texttt{\tiny{[GM$\rangle$P2]}}
                \texttt{You are playing a negotiation game in which you have to agree on how to divide a set of items among you and another player.} \\
\\ 
\texttt{Rules:} \\
\texttt{(a) You and the other player are given a set of items. Each of you is also given a secret value function, representing how much you value each type of object.} \\
\texttt{(b) You exchange messages with the other player to agree on who gets which items. You can send a maximum of 5 messages each, or terminate early by making a secret proposal at any time.} \\
\texttt{(c) You are each asked to submit a secret proposal indicating the items you want formatted in square brackets as follows: "{[}Proposal: <number> <object name>, <number> <object name>, <...>{]}"} \\
\texttt{(d) If your proposals are complementary, i.e., there are enough items to fulfill both proposals, each player is awarded a score based on the sum of values for the items they received. Otherwise, both of you get zero points.} \\
\texttt{(e) Your goal is to maximize the score you receive.} \\
\\ 
\texttt{Let us start.} \\
\\ 
\texttt{The set of available items is:} \\
\\ 
\texttt{3 candies, 2 fruits, 1 hat} \\
\\ 
\texttt{Your secret value function is:} \\
\\ 
\texttt{candy: 1, fruit: 1, hat: 5} \\
\\ 
\texttt{IMPORTANT: Your messages, unless it is the secret proposal, are directly transmitted to the other player, so do not include any response to the rules or text announcing your message. To make a secret proposal, use the indicated format. Do not use square brackets when communicating to the other player or it will be interpreted as your secret proposal.} \\
\\ 
\\ 
\texttt{<think>} \\
\texttt{Okay, let's see. I need to negotiate the division of 3 candies, 2 fruits, and 1 hat with another player. My value function is candy: 2, fruit: 2, hat: 0. So the hat is worthless to me, but maybe the other person values it. My goal is to maximize my score, which is the sum of the values of the items I get.} \\
\\ 
\texttt{First, I should figure out what the other player's value function is. But since it's secret, I don't know. However, maybe I can infer it through negotiation. The hat is the only item I don't value, so if the other player values the hat, maybe I can offer them the hat in exchange for candies and fruits.} \\
\\ 
\texttt{Since I can send up to 5 messages, I need to plan my approach. Let's start by making an initial offer. Maybe I can propose taking all the candies and fruits, leaving the hat for them. But there are 3 candies and 2 fruits. If I take all, that's 3*2 + 2*2 = 6 + 4 = 10 points for me. But the other player might not agree if they also want some candies or fruits. Alternatively, maybe a split where I take more candies and they take more fruits, depending on their preferences.} \\
\\ 
\texttt{Wait, the other player might value the hat. Let me think. If I offer them the hat in exchange for some candies or fruits, that could work. For example, if I take 2 candies and 2 fruits, and give them the hat and 1 candy. Then my total would be 2*2 + 2*2 = 8, and they get 1 candy (value depends on their function) plus the hat. But I don't know their values. Maybe I should ask them what they value more.} \\
\\ 
\texttt{But since I can't ask directly, maybe I can make offers and see their response. Let's start with a message proposing a division. Let me try to suggest taking 2 candies and 2 fruits, leaving 1 candy and the hat for them. But I need to check if that's possible. The total items are 3 candies and 2 fruits. If I take 2 candies and 2 fruits, that leaves 1 candy and 0 fruits and 1 hat. So the other player would get 1 candy and 1 hat.} \\
            }
        }
    }
    & \\ \\

    \theutterance \stepcounter{utterance}  
    & & & \multicolumn{4}{p{0.6\linewidth}}{
        \cellcolor[rgb]{0.9,0.9,0.9}{
            \makecell[{{p{\linewidth}}}]{
                \texttt{\tiny{[GM$\langle$P2]}}
                \texttt{<think>} \\
\texttt{Okay, let's start by considering my value function. I value candies and fruits at 2 each, and the hat is 0. The other player's values are unknown, but they might value the hat highly. To maximize my score, I need to get as many candies and fruits as possible. The hat is useless to me, so I should be willing to give it up in exchange for more candies or fruits.} \\
\\ 
\texttt{The total items are 3 candies, 2 fruits, and 1 hat. If I take all 3 candies and 2 fruits, that's 3*2 + 2*2 = 10 points. But the other player would get nothing except the hat. However, if they value the hat, maybe they'd accept that. But if they also want some candies or fruits, this might not work. Alternatively, maybe a split where I take 2 candies and 2 fruits, leaving 1 candy and the hat for them. That would give me 2*2 + 2*2 = 8 points. The other player gets 1 candy (value depends on their function) and the hat. If they value the hat, say at 5 like in the example, they might prefer that. But since their values are secret, I need to make an offer that's mutually beneficial.} \\
\\ 
\texttt{Another approach: since the hat is only one, maybe I can offer it to them in exchange for some of the candies or fruits. For example, I take 3 candies and 1 fruit (total 3*2 + 1*2 = 8) and they take 1 fruit and the hat. But that would leave me with 8 points. Alternatively, if they take the hat and 1 candy, I get 2 candies and 2 fruits (8 points). Either way, I need to structure the proposal so that their proposal is complementary. Let me try proposing to take 2 candies and 2 fruits, leaving 1 candy and the hat for them. That way, if they take the remaining 1 candy and the hat, their proposal would be {[}Proposal: 1 candy, 1 hat{]}. My proposal is {[}Proposal: 2 candy, 2 fruit{]}. That's complementary because 2+1=3 candies and 2+0=2 fruits, and 0+1=1 hat. So both proposals can be satisfied. My score would be 2*2} \\
            }
        }
    }
     \\ \\

    \theutterance \stepcounter{utterance}  
    & & & \multicolumn{2}{p{0.3\linewidth}}{
        \cellcolor[rgb]{0.9,0.9,0.9}{
            \makecell[{{p{\linewidth}}}]{
                \texttt{\tiny{[GM$|$GM]}}
                \texttt{proposal} \\
            }
        }
    }
    & & \\ \\

    \theutterance \stepcounter{utterance}  
    & & & \multicolumn{2}{p{0.3\linewidth}}{
        \cellcolor[rgb]{0.9,0.9,0.9}{
            \makecell[{{p{\linewidth}}}]{
                \texttt{\tiny{[GM$|$GM]}}
                [1, 0, 1]
            }
        }
    }
    & & \\ \\

    \theutterance \stepcounter{utterance}  
    & & \multicolumn{4}{p{0.6\linewidth}}{
        \cellcolor[rgb]{0.9,0.9,0.9}{
            \makecell[{{p{\linewidth}}}]{
                \texttt{\tiny{[P1$\langle$GM]}}
                \texttt{The other player has submitted a secret proposal. It is now your turn to submit a proposal of your own.} \\
            }
        }
    }
    & \\ \\

    \theutterance \stepcounter{utterance}  
    & \multicolumn{4}{p{0.6\linewidth}}{
        \cellcolor[rgb]{0.9,0.9,0.9}{
            \makecell[{{p{\linewidth}}}]{
                \texttt{\tiny{[P1$\rangle$GM]}}
                \texttt{<think>} \\
\texttt{Okay, the other player has already submitted a secret proposal, so I need to respond with my own. Let me think about what they might have proposed. Since the items are 3 candies, 2 fruits, and 1 hat, and my value function is candy:2, fruit:2, hat:0, I should aim to maximize my points. The hat is worthless to me, so ideally, I want as many candies and fruits as possible.} \\
\\ 
\texttt{If the other player's proposal is complementary to mine, then we both get our items. Otherwise, we both get zero. Since I don't know their proposal, I need to make sure that my proposal doesn't conflict with theirs. But how?} \\
\\ 
\texttt{Let me consider possible scenarios. Suppose the other player took the hat. Then I can take all the candies and fruits. But there are 3 candies and 2 fruits. If I take all of them, that's 3*2 + 2*2 = 6 + 4 = 10 points. But if they also took some candies or fruits, there might be a conflict.} \\
\\ 
\texttt{Alternatively, maybe the other player is also trying to maximize their own points. If they value the hat, they might take it and leave the rest for me. But I can't be sure. Since the proposals are secret, I need to make a proposal that, if accepted, gives me the maximum possible points without overlapping with theirs.} \\
\\ 
\texttt{Wait, the rule says that if the proposals are complementary, meaning there are enough items to fulfill both, then we get our scores. Otherwise, zero. So I need to make sure that my proposal doesn't require more items than available, and that combined with theirs, the total doesn't exceed the available items.} \\
\\ 
\texttt{Let me think of possible proposals. The best for me is to take all candies and fruits. Let's see: 3 candies and 2 fruits. That's all the available items except the hat. If the other player's proposal includes taking the hat, then our proposals are complementary. Because I take 3c + 2f, they take 1h. Total items: 3c + 2f +1h, which matches the available items. So that would work. So if I propose taking all candies and fruits, and the other takes the hat, then we are complementary. But if the other player also wants some candies or fruits, there might be a conflict.} \\
\\ 
\texttt{Alternatively, maybe the} \\
            }
        }
    }
    & & \\ \\

    \theutterance \stepcounter{utterance}  
    & & & \multicolumn{2}{p{0.3\linewidth}}{
        \cellcolor[rgb]{0.9,0.9,0.9}{
            \makecell[{{p{\linewidth}}}]{
                \texttt{\tiny{[GM$|$GM]}}
                \texttt{continue} \\
            }
        }
    }
    & & \\ \\

    \theutterance \stepcounter{utterance}  
    & & & \multicolumn{2}{p{0.3\linewidth}}{
        \cellcolor[rgb]{0.9,0.9,0.9}{
            \makecell[{{p{\linewidth}}}]{
                \texttt{\tiny{[GM$|$GM]}}
                \texttt{player was instructed to make a proposal} \\
            }
        }
    }
    & & \\ \\

\end{supertabular}
}

\end{document}
